% !TeX root = ../../../book.tex
\subsection{示例和伪例}

``语法正确''是指句子中的单词和符号被正确使用和组合并且有意义。这消除了放在一起无意义的符号/单词串，例如：
\begin{center}
    \textcolor{red}{$1+ = 2$, $\text{于儿}^2 = 1$, $\{\{\varnothing\}\} - 7 > 5\pi$, $\text{You am smart}$}
\end{center}
比如，上面的第三个不是数学陈述，因为 $\{\{\varnothing\}\}$ 不是数字，所以我们不知道如何解释从该集合中``减去 $7$''。

而``\emph{唯一}一个真值''说的是该陈述应该要么为\verb|真|要么为\verb|假|，但肯定不能两者兼而有之，或者两者都不是或介于两者之间。这排除了上面``\textcolor{red}{$x^2 - 1 = 0$}''这样的陈述，因为它没有真值。（如果没有声明 $x$ 是什么，我们就无法决定真值是什么。）

\subsubsection*{真值未知}

我们给出的定义中一个奇怪/有趣/复杂的方面是，我们可能不知道给定陈述的真值，即使我们可以确定只有一个这样的值。作为说明，请考虑以下陈述：
\begin{center}
    \textcolor{olivegreen}{任何大于或等于 $4$ 的偶自然数都可以写成两个质数之和。}
\end{center}

这个说法是真是假？如果你能证明或证伪，那么数学界会很乐意看到它！上面的陈述被称为\href{https://baike.baidu.com/item/哥德巴赫猜想/72364}{哥德巴赫猜想}，它是数学中一个非常著名的未解难题（我们希望只是暂时的！）。目前还没有人知道这一说法是对是错，但可以肯定的是，\emph{只有其中一个}真值适用。也就是说，这个陈述不可能既是\verb|真|又是\verb|假|，也不可能介于两者之间。要么所有大于或等于 $4$ 的偶自然数都具有该性质，要么至少有一个不具有该性质。即使还不知道两种可能性中哪一种是正确的，我们也可以陈述这种``非此即彼''的性质。因此，这句话实际上满足了我们对\emph{数学陈述}的定义。

（术语说明：一般来说，\textbf{猜想}是某人认为正确但尚未被证明/证伪的断言。）

\subsubsection*{悖论语句}

使句子没有真值的一种方法是创造\textbf{悖论}。考虑下面这句话：
\begin{center}
    \textcolor{red}{这句话是假的。}
\end{center}
很怪异，对吧？这句话本身就断言了它自己的真值。我们来尝试分析一下它的真值：
\begin{itemize}
    \item 假设这句话为\verb|真|。那么，这个句话本身告诉我们，它实际上为\verb|假|。
    \item 假设这句话为\verb|假|。那么同理，这句话告诉我们，它实际上为\verb|真|。
\end{itemize}
这是不可能的！这句话在某种程度上既\verb|真|又\verb|假|，或者两者都不是，或者……无论是什么，都不是个好主意。我们不想在数学中处理这种奇怪现象，所以我们的定义不允许这类语句作为数学陈述。

（\emph{问题}：如果允许这样的句子成为正确的数学陈述，会发生什么？如果你不遵守我们强调的每个句子必须为\verb|真|或为\verb|假|的原则怎么办？想一想！是不是哪里会出问题，或者这只是一个不同的数学宇宙？……）

通常来说，像上面这样的\textbf{自指}句子（即指代自己的句子）是相当奇怪的，并且可能会产生一些我们想要禁止的悖论。

上述自相矛盾的说法有一个变体，通过下面这幅卡通漫画中给出，其中皮诺曹说：``我的鼻子现在要变长！'' 那么它会变长吗？如果他说真话，那么它就会变长，但只有当他说谎时才它的鼻子才会变长！如果他说谎，那么他的鼻子就会变长（根据定义），但他的说法实际上又是真的！哎呀，纠缠不清！
\begin{center}
    \includegraphics[scale=0.4]{figure/pinocchio.png}
\end{center}

这种现象的一个更奇怪的例子是\emph{奎因悖论(Quine's Paradox)}：
\begin{center}
    \textcolor{red}{``Yields falsehood when preceded by its quotation'' \\yields falsehood when preceded by its quotation.}
\end{center}
这个问题留给你自己思考。我只想说，像这样自相矛盾的说法实在是太病态了，以至于我们不必考虑它们。这就是为什么我们的定义禁止了它们。
